% Reconstructed from the compiled lecture-note PDF.
\chapter{Notation, theorem index, and formula sheet}
This appendix is designed for repeated consultation. The notation follows the source as closely as possible, with occasional simplification explained in the main text.


\section{Geometric objects}

Symbol or term Meaning


K A centrally symmetric convex disk in R2, centered at the origin. Kmin A global minimizer of greatest packing density, equivalently area after normalization. Kccs The class of centrally symmetric convex disks. Kbal The class of balanced disks admitting the six-curve structure used in the control reduction.


δ(K) Greatest packing density of K. For centrally symmetric planar disks, this equals greatest lattice-packing density. L A lattice in R2. det(L) Covolume of the lattice: area of a fundamental parallelogram. HK A least-area centrally symmetric hexagon circumscribing K; a critical hexagon. hK The centrally symmetric hexagon formed by joining six critical boundary points. sj Midpoints of the six edges of a critical hexagon, indexed modulo 6. s∗j The normalized reference multipoint, usually the sixth roots of unity. Tj The equilateral triangular sector containing the jth boundary arc after normalization. σj(t) The six synchronized boundary curves; all have the same image after periodic extension.


Symbol or term Meaning


\begin{displaytext}
multicurve The indexed family (σ0, . . . , σ5) satisfying the multipoint relations at \
every time. \
star condition The tangent σ′j lies in the open cone prescribed by neighboring \
critical points.
\end{displaytext}

smoothed polygon A finite cyclic concatenation of straight segments and tangent hyperbolic arcs of Reinhardt–Mahler type.


\section{Matrix and Lie-theoretic notation}

Symbol Meaning


SL2(R) Real 2 × 2 matrices of determinant 1. sl2(R) Real 2 × 2 matrices of trace 0. J 01 −10 , the infinitesimal rotation. R exp(πJ/3), rotation by π/3. g(t) Matrix curve carrying the reference multipoint to the current one: σj(t) = g(t)s∗j. X(t) Left-trivialized velocity X = g−1g′, normalized by det X = 1. [X, Y ] Matrix commutator XY −Y X. ⟨X, Y ⟩ Trace pairing tr(XY ). Adg X Adjoint action gXg−1. OJ Adjoint orbit of J; the determinant-one sheet containing J. Φ(z) Parametrization of OJ by the upper half-plane. ρ0, ρ1, ρ2 Linear functions of X encoding the star inequalities and edge directions. Zu Control matrix obtained affinely from u = (u0, u1, u2). P Normalized control matrix P = Zu/ ⟨Zu, X⟩.


\section{Control and Hamiltonian notation}

Symbol Meaning


\begin{displaytext}
u = (u0, u1, u2) Normalized curvature control. \
UT Triangular control set ui ≥0, u0 + u1 + u2 = 1. \
e0, e1, e2 Vertices of UT ; bang-bang control values. \
UC, UI, Ur Circular control sets used in the symmetry-enhanced model.
\end{displaytext}

Symbol Meaning


κj State-dependent signed curvature numerator for the even-indexed boundary curves. uj κj/(κ0 + κ1 + κ2). Λ1 Costate paired with the group variable g. Λ2 Original costate paired with X. ΛR Reduced costate [Λ2, X], constrained by ⟨ΛR, X⟩= 0. λ0 Pontryagin cost multiplier; λ0 = −1 in the normal convention. H Control-dependent Hamiltonian. H+ Maximized Hamiltonian maxu∈U H. χij Switching function comparing control vertices i and j. wall A codimension-one set on which two control vertices tie. singular locus The set X = J, ΛR = 0, Λ1 = 32λ0J; geometrically the circle. edge-extremal A trajectory that is also extremal for every restricted edge-control subproblem.


chattering Infinitely many switches accumulating in finite time.


\section{Hyperbolic coordinates}

\begin{displaytext}
For z = x + iy with y > 0, \
x/y −(x2 + y2)/y ! \
Φ(z) = . \
1/y −x/y
\end{displaytext}

The upper-half-plane star domain is


\begin{displaytext}
1 1 h∗= x + iy : −1√ < x < √ x2 + y2 > . 3 3, 3
\end{displaytext}

\begin{displaytext}
The area integrand becomes \
−3 ⟨J, X⟩= 3 x2 + y2 + 1 . \
2 2 y \
The center z = i corresponds to X = J and hence to the circular singular state.
\end{displaytext}

\section{Core equations}

\subsection{Multipoint relations}

\begin{displaytext}
For j modulo 6, \
sj+3 = −sj, \
sj + sj+2 + sj+4 = 0, \
√ \
3 \
det(sj, sj+1) = . \
2
\end{displaytext}

\subsection{Star coordinates}

\begin{displaytext}
For \
a b ! \
X = , \
c −a \
√ \
c + 3a \
ρ0(X) = √ , \
3 \
√ \
c − 3a \
ρ1(X) = √ , \
3 \
+ c ρ2(X) = −3b√ . \
2 3 \
The star condition is ρj(X) > 0 for j = 0, 1, 2, and
\end{displaytext}

\begin{displaytext}
det X = ρ0ρ1 + ρ1ρ2 + ρ2ρ0.
\end{displaytext}

\subsection{Control matrix}

\begin{displaytext}
For u0 + u1 + u2 = 1, √ \
(u1 −u2)/ 3 (u0 −2u1 −2u2)/3 ! \
Zu = √ . \
u0 (u2 −u1)/ 3
\end{displaytext}

It satisfies uj = det(s∗2j, Zus∗2j).


\subsection{State and cost}

\begin{displaytext}
g′ = gX, \
[Zu, X] X′ = ⟨Zu, X⟩, \
Z tf area(K) = −3 ⟨J, X⟩dt. \
2 0 \
Endpoint conditions: \
g(0) = I, g(tf) = R, X(tf) = R−1X(0)R.
\end{displaytext}

\subsection{Hamiltonian system}

\begin{displaytext}
With P = Zu/ ⟨Zu, X⟩, \
Zu⟩ H = Λ1 −3 X −⟨ΛR, 2λ0J, ⟨X, Zu⟩, \
and \
g′ = gX, \
X′ = [P, X], \
Λ′1 = [Λ1, X], \
Λ′R = [P, ΛR] −⟨ΛR, P⟩[P, X] + [−Λ1 + 2λ0J,3 X].
\end{displaytext}

\section{Fuller notation}

Symbol Meaning


w, b, c Complex hyperboloid coordinates centered at the singular locus. z1, z2, z3 Rescaled leading variables of weighted degrees 1, 2, 3. ζ Primitive cube root exp(2πi/3). VT Triangular control values \{1, ζ, ζ2\}. HF Fuller Hamiltonian. AF Fuller angular-momentum analogue. E Exceptional divisor of the weighted blowup. PF Fuller Poincare first-recurrence map on the divisor. PR Exact Reinhardt Poincare map after analytic extension. qin Fixed point representing the inward self-similar chattering trajectory. qout Fixed point representing the outward self-similar trajectory. W s, W u Stable and unstable manifolds of a fixed point.


The triangular Fuller system is


\begin{displaytext}
z′3 = z2, z′2 = z1, z′1 = −iu, u ∈\{1, ζ, ζ2\}.
\end{displaytext}

For constant u and initial values z0j ,


\begin{displaytext}
z1(t) = z01 −itu, \
z2(t) = z02 + z01t −iut2 , 2 \
z01t2 z3(t) = z03 + z02t + −iut3 . 2 6
\end{displaytext}

Therefore every wall comparison is a real cubic in t.


\section{The main theorem chain}

Result Role in the proof


Fejes Toth equality Replaces arbitrary packings by lattice packings for centrally symmetric planar disks. Critical-hexagon theorem Expresses density as area(K)/ area(HK) and supplies six synchronized boundary points.


Balanced-multicurve Converts the minimizer into six regular curves satisfying reduction algebraic and convexity constraints.


Result Role in the proof


\begin{displaytext}
Matrix representation Writes all six curves as σj = gs∗j with g ∈SL2(R). \
State-control theorem Encodes curvature by u ∈UT and obtains \
X′ = [Zu, X]/ ⟨Zu, X⟩. \
Area formula Converts the geometric objective to an integral cost.
\end{displaytext}

Compactification Confines every possible minimizer to a compact subset of the star domain.


Pontryagin maximum Produces costates, switching walls, and necessary principle Hamiltonian dynamics.


Finite switching away from Shows every edge-extremal avoiding the singular locus is a singularity finite smoothed polygon.


Singular-locus classification Identifies full singular arcs with circular arcs.


Second variation Excludes a circular singular interval from a global minimizer.


Fuller truncation Gives the universal leading chattering system near the singular locus.


Weighted blowup Replaces the singular point by the exceptional divisor carrying triangular Fuller dynamics.


Fixed-point theorem Finds the inward and outward self-similar chattering states.


Global basin theorem Forces all nonexceptional divisor trajectories toward the outward fixed point.


Analytic extension Transfers the fixed points and local invariant manifolds to the exact Reinhardt map. Cluster-point theorem Forces every exact finite-time chattering orbit onto W s(qin). No-return theorem Shows the time-reversed unstable curve cannot close into a periodic orbit.


Mahler’s First A periodic global minimizer cannot chatter; hence it is finite bang-bang and a smoothed polygon.


\section{What is analytic and what is computer-assisted}

Category Examples


Classical theorem Fejes Toth, Filippov, Pontryagin, stable-manifold theorem, Weierstrass preparation.


Direct exact calculation Matrix reduction, area formula, state and costate equations, singular-locus classification, Fuller truncation.


Symbolic algebra Switching inequalities, fixed-point polynomials, Jacobians, discriminants, resultants, asymptotic coordinate changes.


Global computer-assisted Fuller cell-image containments and numerical continuation of the step exact unstable curve.


\section{Source-to-notes reading crosswalk}

The notes reorganize rather than duplicate the source. A rough crosswalk is:


Lecture-note material Main source location


Packing geometry and critical Source Chapters 1–2. hexagons


Matrix and control reduction Source Chapters 3–5.


Maximum principle Source Chapter 6.


Bang-bang polygons and Source Chapters 7–8. singularity


Circular model and Fuller Source Chapters 9–12. system


Triangular Fuller map and Source Chapters 13–15. blowup


Global basin and final proof Source Chapters 16 and the concluding part of Part V.


Background structures Source Appendices A–B, supplemented here at a more elementary level.


\section{A one-page proof mnemonic}

Remember the sequence


\begin{displaytext}
packing →critical hexagon →six curves →SL2(R) \
→curvature control →PMP →bang-bang except at the circle
\end{displaytext}

\begin{displaytext}
→chattering model →weighted blowup →global basin
\end{displaytext}

\begin{displaytext}
→unique stable channel →nonperiodic →finite smoothed polygon.
\end{displaytext}

\begin{insightbox}{Chapter summary}
\end{insightbox}

The proof uses many coordinate systems, but each has one job. Planar geometry produces the critical hexagon; SL2(R) synchronizes six curves; the upper half-plane exposes the compact star domain; costates expose switching; complex Fuller coordinates expose chattering; the weighted blowup exposes its limiting directions. The final conclusion is coordinate-independent: the minimizer has only finitely many straight and hyperbolic pieces.

